\section{Đối tượng và phạm vi nghiên cứu}
Đối tượng nghiên cứu của đề tài là nội dung số có khả năng được tạo hoặc chỉnh sửa bởi AI. Trong phạm vi sản phẩm SourceVerify, hệ thống có hướng mở rộng cho ảnh, văn bản và video. Tuy nhiên, do thời lượng Project I có hạn, báo cáo tập trung nhiều nhất vào ảnh vì ảnh số chứa nhiều dấu vết vật lý và thống kê có thể giải thích rõ ràng.

Phạm vi báo cáo gồm:
\begin{itemize}
  \item Trình bày tổng quan bài toán phát hiện nội dung AI.
  \item Phân tích 5 nhóm tín hiệu ảnh nổi bật theo ý tưởng, đầu vào, đầu ra, ưu điểm và hạn chế.
  \item Thiết kế hệ thống SourceVerify ở mức kiến trúc, chức năng và luồng dữ liệu.
  \item Mô tả kết quả triển khai ban đầu của sản phẩm phần mềm.
\end{itemize}
